% ============================================================
% 中文书籍/讲义模板（book-cn）
% ------------------------------------------------------------
% 基于 ctexbook 文档类，默认使用 TeX Live 自带的 fandol 字体，
% 无需在系统中安装额外字体。编译引擎：xelatex。
% 目录结构：
%   main.tex              主文件（扉页 / 前言 / 目录 / 定理环境设置）
%   chapters/chap01.tex   第一章示例
%   chapters/chap02.tex   第二章示例
% 修改指南：
%   1. 在"书籍信息"处填写书名、作者、版本与日期；
%   2. \part 用于分篇，章节数少时可整体删除 \part 行；
%   3. 定理类环境（定义/定理/例）已配好中文名称，直接使用；
%   4. 新增章节：复制 chapters/ 下文件并在 main.tex 中 \include；
%   5. 需要页眉页脚时再引入 fancyhdr，本模板保持默认简洁样式。
% ============================================================
\documentclass[
  UTF8,
  fontset=fandol,        % TeX Live 自带开源字体，跨平台可编译
  zihao=-4,              % 正文字号：小四
  a4paper,
  oneside,               % 单面排版（正式出版可改 twoside）
]{ctexbook}

% ---- 常用宏包 ----
\usepackage{xcolor}                       % 颜色（扉页装饰线等）
\usepackage{amsmath, amssymb}             % 数学公式
\usepackage{amsthm}                       % 定理类环境
\usepackage{graphicx}                     % 插图
\usepackage{booktabs}                     % 三线表
\usepackage{geometry}
\geometry{left=2.8cm, right=2.8cm, top=3cm, bottom=3cm}
\usepackage[
  colorlinks=true,
  linkcolor=black,
  citecolor=black,
  urlcolor=blue,
]{hyperref}

% 数学符号约定：向量用粗斜体
\newcommand{\vect}[1]{\boldsymbol{#1}}

% ---- 定理类环境：中文名称，按章编号 ----
\theoremstyle{definition}
\newtheorem{definition}{定义}[chapter]
\newtheorem{example}{例}[chapter]
\newtheorem{theorem}{定理}[chapter]
\renewcommand{\proofname}{\textbf{证明}}

% ============================================================
% 书籍信息
% ============================================================
\newcommand{\BookTitle}{示例教程}
\newcommand{\BookSubtitle}{从入门到实践}
\newcommand{\BookAuthor}{张~三~编著}
\newcommand{\BookEdition}{第~1~版}
\newcommand{\BookDate}{2026年8月}

\begin{document}

% ==================== 扉页 ====================
\begin{titlepage}
  \begin{center}
    \vspace*{4em}
    {\zihao{1}\heiti \BookTitle}\\[1.2em]
    {\zihao{3} \BookSubtitle}\\[2em]
    {\color{gray}\rule{0.4\textwidth}{0.6pt}}\\[10em]
    {\zihao{3} \BookAuthor}\\[3em]
    {\zihao{3} \BookEdition}\\[0.8em]
    {\zihao{3} \BookDate}
  \end{center}
\end{titlepage}

% ==================== 前言 ====================
\frontmatter  % 前置部分用罗马数字页码（正式页码从第一章开始）
\chapter*{前言}
\addcontentsline{toc}{chapter}{前言}
在此撰写前言。说明本书的写作目的、适用读者、内容组织与
使用建议（例如哪些章节可以跳过、配套资源在哪里获取）。
前言篇幅一般 1～2 页。

% ==================== 目录 ====================
\tableofcontents

% ==================== 正文 ====================
\mainmatter  % 正文页码从 1 重新计数

\part{基础篇}
% 第一章：基础概念（基础篇）
% 示范正文章节的常用元素：小节、定理环境、公式、例题与章末习题。
\chapter{基础概念}

\section{什么是示例问题}
% 开头用一小节引出本章要讲的核心概念
本章介绍全书的基础概念。先用一段话给出直观解释，
再给出严格定义，最后通过例子加深理解。

\begin{definition}[示例问题]
  给定样本集合 $\{(\vect{x}_i, y_i)\}_{i=1}^{N}$，
  示例问题要求寻找参数 $\vect{\theta}$，
  使目标函数 $\mathcal{L}(\vect{\theta})$ 最小。
\end{definition}

\section{基本性质}
% 定理 + 证明的标准组合，证明可留作习题
\begin{theorem}[收敛性]
  当学习率 $\eta$ 满足 $\eta < 1/L$（$L$ 为 Lipschitz 常数）时，
  梯度下降法收敛到局部极小点。
\end{theorem}
\begin{proof}
  由凸性与 Lipschitz 连续性，对目标函数应用标准下降引理即得。
\end{proof}

\section{应用举例}
% 例题环境：先给题目再给解答
\begin{example}[一元情形]
  \label{ex:onedim}
  设 $f(x) = x^2$，从 $x_0 = 2$ 出发，学习率 $\eta = 0.25$，
  求一步梯度下降后的位置。
\end{example}
\begin{proof}[解]
  梯度为 $f'(x) = 2x$，于是
  $x_1 = x_0 - \eta f'(x_0) = 2 - 0.25 \times 4 = 1$。
\end{proof}

\section{本章小结}
% 一段话总结 + 引出下一章
本章给出了示例问题的定义与基本收敛性质，
下一章将在更一般的条件下讨论改进算法。

% ---- 章末习题：讲义常用结构，不需要时可删除 ----
\section*{习\quad 题}
\addcontentsline{toc}{section}{习题}
\begin{enumerate}
  \item 证明：当目标函数为强凸函数时，梯度下降线性收敛。
  \item 编程实现一元梯度下降，并复现例~\ref{ex:onedim} 的计算。
\end{enumerate}


\part{进阶篇}
% 第二章：进阶方法（进阶篇）
% 示范带算法伪代码与表格的章节写法。
\chapter{进阶方法}

\section{改进的优化算法}
% 引出比第一章更一般的算法框架
在基础梯度下降之上，引入动量项可以加速收敛。
更新规则如式~\eqref{eq:momentum}：
\begin{equation}
  \vect{v}_{t+1} = \beta \vect{v}_t - \eta \nabla \mathcal{L}(\vect{\theta}_t),
  \qquad
  \vect{\theta}_{t+1} = \vect{\theta}_t + \vect{v}_{t+1}.
  \label{eq:momentum}
\end{equation}
其中 $\beta \in [0, 1)$ 为动量系数，$\eta$ 为学习率。

\section{超参数的影响}
% 用表格对比不同设置下的表现
表~\ref{tab:hyper} 给出了不同动量系数下的迭代次数对比。

\begin{table}[!htbp]
  \centering
  \caption{不同动量系数下的收敛迭代次数}
  \label{tab:hyper}
  \begin{tabular}{lcccc}
    \toprule
    动量系数 $\beta$ & 0 & 0.5 & 0.9 & 0.99 \\
    \midrule
    收敛迭代次数 & 812 & 356 & 154 & 203 \\
    \bottomrule
  \end{tabular}
\end{table}

\section{本章小结}
动量方法通过累积历史梯度方向改善收敛速度，
适合目标函数存在狭长山谷形等值面的场景；
实际使用中建议从 $\beta = 0.9$ 开始尝试。


% ==================== 附录 ====================
\appendix  % 之后的章节改用字母编号（附录 A、B、…）
\chapter{常用符号一览}
% 附录放辅助性内容：符号表、证明细节、代码清单等
\begin{center}
  \begin{tabular}{ll}
    \toprule
    符号 & 含义 \\
    \midrule
    $\vect{x}$ & 输入向量 \\
    $\vect{\theta}$ & 模型参数 \\
    $\mathcal{L}$ & 损失函数 \\
    \bottomrule
  \end{tabular}
\end{center}

% ==================== 参考文献 ====================
% 教材可用"参考书目"，条目按作者或引用顺序排列
\begin{thebibliography}{9}
\addcontentsline{toc}{chapter}{参考文献}
\bibitem{vapnik2000}
VAPNIK V N. 统计学习理论[M]. 张学工, 译. 北京: 清华大学出版社, 2000.

\bibitem{lecun2015}
LECUN Y, BENGIO Y, HINTON G. Deep learning[J]. Nature, 2015,
521(7553): 436-444.
\end{thebibliography}

\end{document}
